% Part: many-valued-logic
% Chapter: syntax-and-semantics
% Section: semantic-notions

\documentclass[../../../include/open-logic-section]{subfiles}

\begin{document}

\olfileid[id]{mvl}{syn}{sem}

\olsection{Pengertian Semantik}

Andaikan suatu logika bernilai banyak~$\Log L$ diberikan oleh sebuah matriks.
Kita kemudian dapat mendefinisikan pengertian semantik biasa bagi~$\Log L$.

\begin{defn}
\begin{enumerate}
\item !!^a{formula}~$!A$ \emph{dapat dipenuhi} jika terdapat
  $\pAssign{v}$ sedemikian sehingga $\pSat{v}{!A}$; formula itu
  \emph{tidak dapat dipenuhi} jika tidak ada $\pAssign{v}$ sedemikian
  sehingga $\pSat{v}{!A}$;
\item !!^a{formula}~$!A$ merupakan \emph{tautologi} jika
  $\pSat{v}{!A}$ bagi semua !!{valuation}s~$v$;
\item Jika $\Gamma$ merupakan suatu himpunan !!{formula}s, maka
  $\Gamma\Entails !A$ (``$\Gamma$ mengikutkan~$!A$'') jika dan hanya jika
  $\pSat{v}{!A}$ bagi setiap !!{valuation}~$\pAssign{v}$ yang memenuhi
  $\pSat{v}{\Gamma}$.
\item Jika $\Gamma$ merupakan suatu himpunan !!{formula}s, maka $\Gamma$
  \emph{dapat dipenuhi} jika terdapat !!a{valuation}~$\pAssign{v}$ yang
  memenuhi $\pSat{v}{\Gamma}$, dan $\Gamma$ \emph{tidak dapat dipenuhi}
  dalam kasus lainnya.
\end{enumerate}
\end{defn}

Beberapa fakta mengenai pengertian-pengertian ini sama seperti dalam kasus
logika klasik:

\begin{prop}
\ollabel{prop:semanticalfacts}
\begin{enumerate}
\item $!A$ merupakan tautologi jika dan hanya jika
  $\emptyset\Entails !A$;
\item Jika $\Gamma$ dapat dipenuhi, maka setiap himpunan bagian berhingga
  dari~$\Gamma$ juga dapat dipenuhi;
\item\ollabel{def:monotonicity}%
Monotonisitas: jika $\Gamma\subseteq\Delta$ dan $\Gamma\Entails !A$, maka
  $\Delta\Entails !A$;
\item\ollabel{def:Cut}%
Transitivitas: jika $\Gamma\Entails !A$ dan
  $\Delta\cup\{!A\}\Entails !B$, maka
  $\Gamma\cup\Delta\Entails !B$.
\end{enumerate}
\end{prop}

\begin{proof}
Latihan.
\end{proof}

\begin{prob}
Buktikan \olref[mvl][syn][sem]{prop:semanticalfacts}.
\end{prob}

Dalam logika klasik kita dapat menghubungkan perikutan dengan implikasi.
Sebagai contoh, kita mempunyai validitas \emph{modus ponens}: jika
$\Gamma\Entails !A$ dan $\Gamma\Entails !A\lif !B$, maka
$\Gamma\Entails !B$. Hubungan penting lainnya antara $\Entails$ dan~$\lif$
dalam logika klasik ialah teorema deduksi semantik:
$\Gamma\Entails !A\lif !B$ jika dan hanya jika
$\Gamma\cup\{!A\}\Entails !B$. Hasil-hasil ini \emph{tidak} selalu berlaku
dalam logika bernilai banyak. Berlaku atau tidaknya hasil-hasil tersebut
bergantung pada fungsi kebenaran~$\tf{\lif}$.

\end{document}
